\documentclass[11pt]{article}
\usepackage[T1]{fontenc}
\usepackage{lmodern}
\usepackage[margin=1in]{geometry}
\usepackage{amsmath,amssymb,amsthm,booktabs,array}
\usepackage[colorlinks=true,urlcolor=blue,linkcolor=blue,citecolor=blue]{hyperref}
\usepackage{microtype}
\newtheorem{theorem}{Theorem}
\newtheorem{lemma}[theorem]{Lemma}
\newtheorem{proposition}[theorem]{Proposition}
\newcommand{\Z}{\mathbb Z}
\newcommand{\Q}{\mathbb Q}
\newcommand{\OK}{\mathcal O_K}
\newcommand{\zK}{\zeta_K}
\newcommand{\tht}{\Theta}
\newcommand{\norm}[1]{\lvert #1\rvert}
\title{An explicit tenth-discriminant bound\\from Banks's zero-free theorem}
\author{Research note prepared with Codex}
\date{20 October 2026}
\begin{document}
\maketitle
\begin{abstract}
Taking the explicit large-conductor theorem of Banks
(\texttt{arXiv:2610.15483v1}, Theorem~1.1) as an external input, we show that
an imaginary quadratic field of class number one has absolute fundamental
discriminant less than \(10^{10}\).
The additional analytic argument is a short, explicit application of
Poisson summation to the integer lattice of the field.
An exhaustive integer calculation using only primes at most \(109\)
then recovers exactly the nine familiar fundamental discriminants.
Two implementations, using different algorithms, give the same result.
The argument does not use Heegner's solution, Heegner points, or a
precomputed class-number table.
\end{abstract}

\section{Statement and dependency}

Throughout, \(D<0\) is a \emph{fundamental discriminant},
\(q=\norm{D}\), \(K=\Q(\sqrt D)\), and \(\chi_D\) is the primitive
quadratic character of conductor \(q\). All logarithms are natural.
In particular, \(q\) is the absolute field discriminant, not necessarily
the squarefree integer appearing under a square root.

The only recent theorem used is the following assertion of Banks~\cite{banks}:
\begin{equation}\label{banks}
 q\ge 10^{10},\quad L(\beta,\chi)=0,\quad 0<\beta<1
 \quad\Longrightarrow\quad
 (1-\beta)\log q\ge \frac{3}{1000}
\end{equation}
for every primitive real character \(\chi\) of conductor \(q\).
Only this large-conductor assertion is needed. Its extension to smaller
conductors by a separate computation is not used.

\begin{theorem}\label{main}
Assuming the external theorem \eqref{banks}, if \(h(K)=1\), then
\[
 q<10^{10}.
\]
Combining this bound with the finite verification described in
Section~\ref{finite} gives
\[
 q\in\{3,4,7,8,11,19,43,67,163\}.
\]
Every discriminant on this list does have class number one.
\end{theorem}

The qualifier in the statement records the precise external dependency:
this note does not supply an independent proof of Banks's theorem.
The contribution here is the explicit passage from that theorem to the
class-number-one problem, together with a reproducible finite exclusion.
No novelty or optimality claim is made for the underlying classical
theta-function argument.

\section{Class number one produces a nearby real zero}

\begin{lemma}\label{zero}
If \(q\ge10^{10}\) and \(h(K)=1\), then \(L(s,\chi_D)\) has a real zero
\[
 1-\frac{10}{\sqrt q}<\beta<1.
\]
\end{lemma}

\begin{proof}
Put
\[
 R=\sqrt q,\qquad y=\frac R2,\qquad
 \Lambda=y^{-1/2}\OK\subset\mathbb C\simeq\mathbb R^2.
\]
The embedded integer lattice \(\OK\) has covolume \(y\), so
\(\Lambda\) has covolume one. Define
\[
 \tht(t)=\sum_{\alpha\in\OK}
       \exp\left(-\frac{\pi t\norm{\alpha}^2}{y}\right)
       \qquad(t>0).
\]
For any planar lattice of covolume one, its Euclidean dual is a
rotation by \(90\) degrees of the original lattice. For example, if
the columns of \(A\) are an oriented basis and \(\det A=1\), then
\(A^{-T}=JAJ^{-1}\), where \(J\) is rotation by \(90\) degrees
and \(J^{-1}\Z^2=\Z^2\).
Consequently, Poisson summation for the Gaussian gives
\begin{equation}\label{thetafunctional}
 \tht(t)=t^{-1}\tht(1/t).
\end{equation}

Since \(q>4\), the units of \(\OK\) are exactly \(\{\pm1\}\).
Class number one implies that every nonzero integral ideal has
exactly two generators in \(\OK\). Thus, initially for \(\Re s>1\),
\[
 \sum_{\alpha\in\OK\setminus\{0\}}\norm{\alpha}^{-2s}
     =2\zK(s).
\]
Taking the Mellin transform and splitting its integral at \(t=1\),
then using \eqref{thetafunctional} in the part below \(1\), gives
\begin{equation}\label{mellin}
 2\Gamma(s)\left(\frac y\pi\right)^s\zK(s)
 =
 -\frac{1}{s(1-s)}
 +\int_1^\infty(\tht(t)-1)
       \left(t^{s-1}+t^{-s}\right)\,dt .
\end{equation}
The integral on the right converges locally uniformly for all complex
\(s\), since \(\tht(t)-1\) decays exponentially as \(t\to\infty\).
Hence \eqref{mellin} continues the initial identity meromorphically.
For clarity, the rational term is
\[
 \int_0^1(t^{s-2}-t^{s-1})\,dt
 =\frac{1}{s-1}-\frac1s=-\frac{1}{s(1-s)}.
\]

Set
\[
 \varepsilon=\frac{10}{R},\qquad s_0=1-\varepsilon,\qquad
 A=\frac y\pi=\frac{R}{2\pi}.
\]
Here \(R\ge10^5\), \(0<s_0<1\), and \(A>1\).
The two lattice points \(\alpha=1,-1\) alone give
\[
 \tht(t)-1\ge2e^{-\pi t/y}\ge\frac2e
       \qquad(1\le t\le A).
\]
Dropping the nonnegative \(t^{-s_0}\) term and the part of the integral
above \(A\), the integral in \eqref{mellin} is therefore at least
\[
 \frac2e\int_1^A t^{s_0-1}\,dt
    =\frac{2}{e\,s_0}(A^{s_0}-1).
\]
It suffices to prove
\begin{equation}\label{target}
 \frac{2\varepsilon}{e}(A^{1-\varepsilon}-1)>1.
\end{equation}

All estimates needed for \eqref{target} can be made rational.
The function \((\log R)/R\) decreases for \(R>e\), and
\(\log(10^5)<12\). Hence
\[
 0<\varepsilon\log A
 \le\frac{10\log R}{R}<\frac{3}{2500}.
\]
Using \(e^{-u}\ge1-u\),
\[
 A^{1-\varepsilon}
   = A e^{-\varepsilon\log A}
   >\frac{R}{2\pi}\frac{2497}{2500}.
\]
The elementary bounds \(\pi<22/7\) and \(2<e<3\) now imply
\begin{align*}
 \frac{2\varepsilon}{e}(A^{1-\varepsilon}-1)
 &>\frac{10}{\pi e}\frac{2497}{2500}-\frac{20}{eR}\\
 &>\frac{35}{33}\frac{2497}{2500}-\frac{1}{10000}\\
 &=\frac{31777}{30000}>1.
\end{align*}
For an entirely rational check of the logarithmic bound, use
\(e>8/3\) and \(8^{12}>10^5\,3^{12}\).
This proves \eqref{target}, and so \(\zK(s_0)>0\).

As \(s\to1^-\), the rational term in \eqref{mellin} tends to
\(-\infty\), while its integral stays finite. Thus \(\zK(s)<0\)
sufficiently close to \(1\) from the left.
Continuity gives a zero \(\beta\in(s_0,1)\).
Finally \(\zK(s)=\zeta(s)L(s,\chi_D)\), and \(\zeta(s)\ne0\) for
real \(0<s<1\). One elementary proof of the latter uses the positive
alternating series
\[
 \eta(s)=\sum_{n\ge1}(-1)^{n-1}n^{-s}>0,\qquad
 \zeta(s)=\frac{\eta(s)}{1-2^{1-s}}<0.
\]
Therefore \(L(\beta,\chi_D)=0\), as claimed.
\end{proof}

\begin{proof}[Proof of the analytic bound in Theorem~\ref{main}]
If \(q\ge10^{10}\) and \(h(K)=1\), Lemma~\ref{zero} gives
\[
 (1-\beta)\log q < \frac{10\log q}{\sqrt q}.
\]
The function \((\log q)/\sqrt q\) decreases for \(q>e^2\).
Using \(\log(10^{10})<24\), we obtain
\[
 (1-\beta)\log q
 < \frac{10\cdot24}{10^5}
 =\frac{3}{1250}<\frac{3}{1000}.
\]
This contradicts \eqref{banks}.
\end{proof}

\section{A finite exclusion with small primes}\label{finite}

\begin{lemma}\label{inert}
If \(h(K)=1\), then every prime \(p<q/4\) is inert in \(K\).
\end{lemma}
\begin{proof}
If \(p\) splits or ramifies, there is an integral ideal of norm \(p\).
It is principal when \(h(K)=1\), so it has a generator
\(\alpha\in\OK\) with \(\norm{\alpha}^2=p\).
A rational integer cannot have prime norm in a quadratic field.
Every nonrational algebraic integer in \(K\) can be written
\((a+b\sqrt{-q})/2\), with integers \(a,b\) and \(b\ne0\);
the allowed parities depend on \(D\).
Consequently,
\[
 p=\norm{\alpha}^2=\frac{a^2+qb^2}{4}\ge\frac q4,
\]
a contradiction.
\end{proof}

For \(q>8\), Lemma~\ref{inert} applies to \(p=2\).
Thus \(2\) is inert, \(D\) is odd, and \(D\equiv5\pmod8\);
equivalently, \(q\equiv3\pmod8\).
The negative fundamental discriminants with \(q\le8\) are
\(-3,-4,-7,-8\).

For every remaining candidate, therefore, it is enough to examine
\[
 11\le q<10^{10},\qquad q\equiv3\pmod8.
\]
For an odd prime \(p<q/4\), exclude \(q\) whenever \(-q\) is zero
or a square modulo \(p\).
The search deliberately includes nonfundamental \(q\), so no
squarefreeness test can accidentally discard a valid candidate.

\subsection*{Exact computation}
Two programs accompany this note:
\begin{itemize}
\item \texttt{verify\_sieve.cpp} builds permitted residue classes using
the Chinese remainder theorem, then checks candidates individually
using tables of squares modulo primes.
\item \texttt{verify\_independent.py} obtains the initial residue
classes by direct scanning and Euler's criterion. It then checks
whole arithmetic progressions using periodic integer bitsets.
\end{itemize}

Both first impose the conditions for \(2,3,5,7,11,13,17\).
The resulting wheel has modulus
\[
 M=8\cdot3\cdot5\cdot7\cdot11\cdot13\cdot17
   =2042040
\]
and \(1440\) permitted residue classes.
For \(437\le q<10^{10}\), every prime at most \(109\) satisfies
\(4p<q\). The programs check the remaining primes from \(19\)
through \(109\), and find no surviving \(q\).
For \(11\le q<437\), they apply only the prime conditions
with \(4p<q\), and find exactly
\[
 11,\quad19,\quad43,\quad67,\quad163.
\]
The endpoint \(10^{10}\) is excluded analytically, and is the
exclusive upper endpoint in both programs.
These calculations use exact integer arithmetic throughout.

\subsection*{Checking the surviving fields}
The Python verifier independently enumerates all primitive reduced
positive binary quadratic forms of each remaining discriminant.
The unique reduced form in each case is as follows:
\[
\begin{array}{c|rrrrrrrrr}
q&3&4&7&8&11&19&43&67&163\\ \hline
(a,b,c)&
(1,1,1)&(1,0,1)&(1,1,2)&(1,0,2)&
(1,1,3)&(1,1,5)&(1,1,11)&(1,1,17)&(1,1,41)
\end{array}
\]
In the enumeration, \(\norm b\le a\le c\),
\(b^2-4ac=-q\), \(\gcd(a,b,c)=1\), and \(b\ge0\) on
the boundaries \(\norm b=a\) or \(a=c\).
The standard bound \(a\le\sqrt{q/3}\) makes the enumeration finite.
The classical correspondence with ideal classes verifies
class number one for these nine fields.

\section{Reproduction, scope, and audit}

Run the following commands from the directory containing the sources:
\begin{verbatim}
g++ -O2 -std=c++17 -Wall -Wextra -pedantic \
    verify_sieve.cpp -o verify_sieve
./verify_sieve
python3 verify_independent.py
\end{verbatim}
The programs require no number-theory library or external data.
Their outputs are saved in \texttt{sieve\_cpp\_output.txt} and
\texttt{sieve\_python\_output.txt}.
The Python program also checks the exact rational inequalities in
the analytic argument.

The source snapshot and its SHA-256 hash are recorded in
\texttt{sources/manifest.json}. The recent external dependency is
Banks's stated Theorem~1.1 for \(q\ge10^{10}\); this note
does not certify its full nineteen-page proof.
Classical inputs are Poisson summation for a Gaussian,
basic quadratic ideal theory and the factorization
\(\zK=\zeta L\), and reduction of positive binary quadratic forms.
The analytic lemma is proved above rather than imported with an
unspecified error constant.

The constant \(10\) in Lemma~\ref{zero} is chosen for a short proof
with a comfortable margin. It is unrelated to the illustrative
additive constant \(C_{\rm LS}=10\) in the blog comments.
Sharpening this lemma alone would not lower the large-conductor
threshold in the quoted form of Banks's theorem.
The finite exclusion already gives the exact final classification.

This is a candidate contribution to Tao's challenge~\cite{tao},
with a clear external theorem dependency and reproducible arithmetic.
It is not a claim to have independently established the new
zero-free theorem, and no Lean formalization has been produced.

\paragraph{AI disclosure.}
The route selection, explicit estimates, proof text, and both verification
programs were produced by Codex in this conversation.
The two computations are algorithmically independent, but were written
by the same AI assistant; they are not an independent human or
independent-model review.
The user has not yet reviewed or endorsed the mathematical argument.
Nothing has been submitted to the blog or published externally.

\begin{thebibliography}{9}
\bibitem{banks}
W.~D.~Banks,
\emph{On the OpenAI proof that Landau--Siegel zeros do not exist},
11 October 2026, arXiv:2610.15483v1, Theorem~1.1.
\url{https://arxiv.org/html/2610.15483v1}.
\bibitem{tao}
T.~Tao,
\emph{An experimental AI-math challenge: using tenth discriminant
bounds to explore arguments ruling out Siegel zeroes},
20 October 2026.
\href{https://terrytao.wordpress.com/2026/10/20/an-experimental-ai-math-challenge-using-tenth-discriminant-bounds-to-explore-arguments-ruling-out-siegel-zeroes/}
{Online challenge}.
\end{thebibliography}
\end{document}
